\documentclass[conference]{IEEEtran}
\usepackage{graphicx}
\usepackage{booktabs}
\usepackage{amsmath}
\usepackage{cite}
\graphicspath{{../results/figures/}{results/figures/}}

\begin{document}

\title{A Lightweight Tabular Machine-Learning Pipeline for HER/OER Catalyst Screening}

\author{\IEEEauthorblockN{Course Project Draft}
\IEEEauthorblockA{Graduate Machine Learning Term Project}}

\maketitle

\begin{abstract}
This project implements a reproducible, CPU-only machine-learning workflow for screening electrocatalyst candidates for hydrogen evolution reaction (HER) and oxygen evolution reaction (OER) descriptors. The goal is not to claim new state-of-the-art catalysts, but to compare how two electrochemical reactions behave under one consistent tabular modeling framework. Reaction energies were queried from Catalysis-Hub with bounded pagination and cached to disk. The final run used only real Catalysis-Hub data: 1200 HER raw rows and 12362 OER adsorbate rows before cleaning and grouping. Features combine matminer Magpie composition descriptors with transparent hand-crafted elemental and site descriptors. Linear, ridge, random-forest, XGBoost, and small MLP regressors were evaluated using five-fold cross-validation. The best HER model was a random forest with MAE 1.108 eV and $R^2=0.373$. Mean best OER-target performance was MAE 0.539 and $R^2=0.561$ across the four OER targets. The pipeline outputs metrics, feature importances, parity plots, volcano plots, rankings, and a comparison summary.
\end{abstract}

\begin{IEEEkeywords}
electrocatalysis, HER, OER, Catalysis-Hub, matminer, random forest, XGBoost, volcano plot
\end{IEEEkeywords}

\section{Introduction and Problem Definition}
Water splitting requires efficient electrocatalysts for both hydrogen evolution reaction (HER) and oxygen evolution reaction (OER). A standard computational screening strategy is to model adsorption free-energy descriptors that control activity. HER is commonly summarized by the hydrogen adsorption free energy, $\Delta G_{H*}$, with an ideal volcano optimum near 0 eV. OER is more complex because it involves a four-step pathway through OH*, O*, and OOH* intermediates and is often summarized by a theoretical overpotential, $\eta$.

This project builds a complete graduate-course-scale machine-learning pipeline for predicting these descriptors from tabular elemental, compositional, and limited site descriptors. The pipeline screens candidates by ranking HER materials toward $|\Delta G_{H*}|=0$ and OER materials toward low predicted overpotential. The main comparison question is whether HER and OER are equally predictable under the same features and model classes, and whether top candidates overlap enough to motivate bifunctional catalyst discussion.

\section{Related Work}
Adsorption-energy descriptors and volcano plots are widely used in computational catalysis because they reduce complex surface kinetics to interpretable thermodynamic proxies. The Sabatier principle motivates the HER target: hydrogen should bind neither too strongly nor too weakly. For OER, four proton-electron transfer steps define the thermodynamic limiting step, and the overpotential follows from the largest step free energy relative to 1.23 V.

Materials informatics libraries such as pymatgen and matminer provide reproducible composition-based descriptors. Tree ensembles and gradient boosting are common baselines for tabular materials data because they handle nonlinear elemental trends without requiring large datasets. Neural networks are included here only as a small comparison model, not as a deep-learning architecture.

\section{Methodology}
\subsection{Data Acquisition}
The acquisition module queries the Catalysis-Hub GraphQL endpoint with a small probe query before extraction. The implemented extractor uses bounded cursor pagination, parses product JSON keys locally, and filters exact adsorbate products Hstar, OHstar, Ostar, and OOHstar. All network data are cached to \texttt{data/raw/} and reruns use the cache. This keeps the workflow offline after the first successful pull and avoids large archives such as OC20.

The HER data path produced 1200 raw Catalysis-Hub rows. To avoid synthetic OER data, the OER scan cap was increased within the project constraint to 50000 bounded reaction records. This produced 12362 real OER adsorbate rows and 94 complete candidate groups after pivoting OH*, O*, and OOH* into one row per candidate.

The exact extraction logic was kept deliberately simple because the GraphQL schema exposes reaction fields but adsorbate filtering can be deployment-dependent. Each returned reaction node contains a formula, surface composition, facet, product dictionary, reaction string, site dictionary, and reaction energy. Product keys are mapped only when they exactly match the intended adsorbates, avoiding accidental matches such as larger carbon-containing adsorbates. The cache files therefore preserve both the original reaction metadata and the derived target labels used by the ML pipeline.

\subsection{Data Completeness Policy}
The code still contains a documented fallback path for robustness, but it was not used in the final reported run. The final raw and processed OER files have \texttt{source=Catalysis-Hub} and \texttt{fallback\_used=False}. The main tradeoff is that real complete OER triplets are much fewer than single-adsorbate rows, so the OER model is trained on a smaller but more internally consistent subset.

\subsection{Featurization}
Formulas are parsed with pymatgen where available and with a simple internal parser otherwise. The primary descriptors are matminer \texttt{ElementProperty.from\_preset("magpie")} features. Additional hand-crafted features include weighted means, standard deviations, minima, and maxima of atomic number, electronegativity, atomic radius, group, period, valence-electron proxy, d-electron proxy, oxidation-state proxy, oxygen fraction, site-type indicators, a coordination-number estimate, and a simple d-band proxy. All-NaN and constant columns are removed and remaining missing values are median-imputed.

For OER, long-form adsorbate rows are pivoted into one row per candidate. The overpotential is computed from
\begin{align}
\Delta G_1 &= \Delta G_{OH*},\\
\Delta G_2 &= \Delta G_{O*} - \Delta G_{OH*},\\
\Delta G_3 &= \Delta G_{OOH*} - \Delta G_{O*},\\
\Delta G_4 &= 4.92 - \Delta G_{OOH*},\\
\eta &= \max_i(\Delta G_i) - 1.23.
\end{align}
Because each step transfers one electron, the eV value is numerically equivalent to volts for this descriptor calculation.

The feature matrix is intentionally tabular and CPU-friendly. It does not use graph neural networks, slab images, or large-scale pretraining. This makes the experiment appropriate for a one-semester ML course: the focus is on acquisition, cleaning, feature engineering, cross-validation, and interpretation rather than on heavy representation learning.

\section{Model Architecture}
Five model families are evaluated for HER and for each OER target: ordinary least squares linear regression, ridge regression, random forest regression, XGBoost regression, and a small MLP regressor. Linear, ridge, and MLP models use standard scaling. Random forest and XGBoost use small CPU-friendly hyperparameter grids read from \texttt{config.yaml}. The MLP contains two hidden layers of sizes 64 and 32 with early stopping.

For computational control, the workflow fixes random seed 42, uses five-fold shuffled cross-validation, limits model samples per target, and selects the top-variance numeric features for modeling. The selected feature cap preserves the course-project requirement to stay CPU-only while still using matminer descriptors.

The random forest and XGBoost grids are intentionally small. Random forest varies tree depth while using a fixed modest number of estimators; XGBoost varies learning rate with shallow histogram trees. This is not an exhaustive hyperparameter optimization. It is a light search designed to demonstrate the methodology while keeping runtime practical on a standard CPU. The best model is selected by MAE because adsorption-energy ranking is usually more sensitive to absolute energy errors than to explained variance alone.

\section{Experimental Setup}
The final processed matrices contained 937 HER rows with 178 columns and 94 OER rows with 158 columns. Metrics are mean absolute error (MAE), root mean squared error (RMSE), and coefficient of determination ($R^2$), reported as mean and standard deviation over five folds. The best model per target is selected by lowest cross-validated MAE, then refit to generate candidate predictions and parity plots.

All generated files are organized to support auditability. Raw data are in \texttt{data/raw/}, processed matrices in \texttt{data/processed/}, figures in \texttt{results/figures/}, metrics and importances in \texttt{results/metrics/}, rankings in \texttt{results/rankings/}, and written summaries in \texttt{results/} and \texttt{report/}. The command \texttt{python run\_all.py} reruns the full workflow from cache and reproduces the reported numbers in this draft.

\begin{table}[t]
\centering
\caption{Best cross-validated model per target. Numeric values are generated by \texttt{results/metrics/cv\_scores.csv}.}
\begin{tabular}{llccc}
\toprule
Reaction & Target & Model & MAE & $R^2$ \\
\midrule
HER & $\Delta G_{H*}$ & RF & $1.108\pm0.086$ & $0.373\pm0.046$ \\
OER & $\Delta G_{OH*}$ & XGB & $0.465\pm0.124$ & $0.765\pm0.132$ \\
OER & $\Delta G_{O*}$ & XGB & $0.659\pm0.132$ & $0.703\pm0.085$ \\
OER & $\Delta G_{OOH*}$ & XGB & $0.576\pm0.118$ & $0.663\pm0.108$ \\
OER & $\eta$ & XGB & $0.456\pm0.176$ & $0.112\pm0.954$ \\
\bottomrule
\end{tabular}
\end{table}

\section{Results and Analysis}
\subsection{EDA}
Figure~\ref{fig:eda} shows representative target distributions and feature-space structure. The PCA plots are used only as descriptive visualizations; model selection is based on cross-validation rather than visual separation.

\begin{figure}[t]
\centering
\includegraphics[width=0.48\linewidth]{her_target_histograms.png}
\includegraphics[width=0.48\linewidth]{oer_target_histograms.png}
\includegraphics[width=0.48\linewidth]{her_dG_H_pca.png}
\includegraphics[width=0.48\linewidth]{oer_overpotential_pca.png}
\caption{Target distributions and two-dimensional PCA summaries for HER and OER feature matrices.}
\label{fig:eda}
\end{figure}

Correlation heatmaps provide a sanity check that the model has access to chemically meaningful variation. Figure~\ref{fig:corr} shows that HER and OER do not emphasize identical feature-target relationships. This supports the later importance analysis: even when the same descriptor generator is used, each reaction selects different slices of the composition space.

\begin{figure}[t]
\centering
\includegraphics[width=0.48\linewidth]{her_dG_H_correlation_heatmap.png}
\includegraphics[width=0.48\linewidth]{oer_overpotential_correlation_heatmap.png}
\caption{Top feature-target correlation heatmaps for HER $\Delta G_{H*}$ and OER overpotential.}
\label{fig:corr}
\end{figure}

\subsection{Prediction Performance}
The best HER model was random forest with MAE 1.108 eV and $R^2=0.373$. XGBoost was similar but slightly worse by MAE. HER therefore showed modest predictability from composition and simple site descriptors. OER targets were best predicted by XGBoost in this real-data run, with strong $R^2$ for the three adsorption energies and weaker, more variable $R^2$ for the derived overpotential. This is expected because overpotential is a max-step transformation and therefore amplifies errors in the intermediate predictions.

\begin{figure}[t]
\centering
\includegraphics[width=0.48\linewidth]{parity_her_dG_H.png}
\includegraphics[width=0.48\linewidth]{parity_oer_overpotential.png}
\caption{Parity plots for the selected HER descriptor model and OER overpotential model.}
\label{fig:parity}
\end{figure}

The comparison summary concludes that OER was easier by average $R^2$ in this bounded real-data run. This should be interpreted carefully: complete OER triplets come from a smaller and more internally consistent subset of Catalysis-Hub, whereas the HER cache spans broader chemistry and duplicated site environments.

A second observation is the role of linear baselines. Ridge regression was competitive for OER $\Delta G_{O*}$, while tree models were best for HER, OER $\Delta G_{OH*}$, and OER overpotential. This outcome is useful pedagogically: simple regularized models remain important references, and a nonlinear model should be interpreted relative to those baselines rather than reported in isolation. The MLP was not competitive in this small-data setting, which is consistent with the expectation that neural networks require either more data or stronger representation learning to outperform tree ensembles on tabular materials descriptors.

\subsection{Feature Importance}
Tree-model importance files show that HER relies heavily on total atom count and radius statistics, followed by Magpie ground-state volume and melting-temperature descriptors. OER feature importance emphasizes Magpie melting-temperature deviation, radius standard deviation, covalent-radius statistics, total atoms, and space-group-number derived composition statistics. These differences suggest that the two reactions are sensitive to overlapping but not identical compositional trends.

Feature importance should not be interpreted causally. The random forest and XGBoost importances are impurity/gain-style model diagnostics, so correlated descriptors can share or exchange importance. They are best used here as an interpretability guide: they identify which compositional statistics the fitted model used most strongly, but they do not prove a mechanistic adsorption-energy relationship.

\subsection{Screening and Volcano Plots}
The HER screening score is $|\widehat{\Delta G}_{H*}|$, ranked ascending after aggregating duplicate rows to unique candidates. The top predicted HER candidate was \texttt{Ag9Ru3|Ag3Ru|111}, with predicted $\Delta G_{H*}=-0.005$ eV. The OER screening score is predicted overpotential, ranked ascending. The top OER candidate was \texttt{IrC26N4|IrN4C26-Delta z=relaxed-M|001}, with predicted $\eta=0.605$ V. Both rankings in the final run are based on real Catalysis-Hub rows.

\begin{figure}[t]
\centering
\includegraphics[width=0.48\linewidth]{her_volcano.png}
\includegraphics[width=0.48\linewidth]{oer_volcano.png}
\caption{Volcano-style screening plots. HER is optimized near predicted $\Delta G_{H*}=0$; OER is ranked by low predicted overpotential versus the predicted $\Delta G_{O*}-\Delta G_{OH*}$ descriptor.}
\label{fig:volcano}
\end{figure}

There was no exact top-25 candidate overlap and no top-25 formula overlap between HER and OER. The bifunctional implication is that single-objective HER and OER rankings should not be merged naively. A future bifunctional search should explicitly optimize both scores, for example using a Pareto front.

Both top lists contain Catalysis-Hub candidates in the final run. There is still no top-25 overlap, so the qualitative conclusion is that HER and OER objectives select different regions of this bounded candidate space.

\begin{figure}[t]
\centering
\includegraphics[width=0.92\linewidth]{her_oer_comparison_metrics.png}
\caption{Best-model MAE and $R^2$ comparison across HER and OER targets.}
\label{fig:comparison}
\end{figure}

\section{Limitations and Future Work}
The main limitation is data completeness. A bounded Catalysis-Hub pull was sufficient for HER but not for complete OER OH/O/OOH groups, so OER results are limited by the documented fallback dataset. This is appropriate for a reproducible course-project pipeline but not sufficient for scientific catalyst discovery. A second limitation is feature scope: composition descriptors ignore explicit slab geometry, adsorption site identity, coverage, and electronic structure. Finally, random train/test splits may overestimate generalization for chemically similar compositions.

Future work should implement adsorbate-specific Catalysis-Hub queries if a stable API filter becomes available, add explicit surface graph or local coordination descriptors, use group splits by composition family, calibrate uncertainty for screening, and run a true multi-objective HER/OER bifunctional optimization.

Additional validation should include repeated cross-validation, leave-composition-family-out splits, and calibration of prediction intervals. Screening based only on point predictions can over-rank uncertain candidates. A practical next version would report both expected activity and uncertainty, then prioritize candidates that are simultaneously promising and reliable.

\section*{Reproducibility Statement}
All random seeds are centralized in \texttt{config.yaml}. Raw network pulls are cached in \texttt{data/raw/}. Processed matrices, metrics, feature importances, rankings, and figures are generated by \texttt{python run\_all.py}. The workflow is CPU-only and does not download large bulk archives.

\begin{thebibliography}{1}
\bibitem{catalysishub} S. Winther et al., ``Catalysis-Hub.org, an open electronic structure database for surface reactions,'' \emph{Scientific Data}, vol. 6, 2019.
\bibitem{matminer} L. Ward et al., ``Matminer: An open source toolkit for materials data mining,'' \emph{Computational Materials Science}, vol. 152, pp. 60--69, 2018.
\bibitem{pymatgen} S. P. Ong et al., ``Python Materials Genomics (pymatgen): A robust, open-source python library for materials analysis,'' \emph{Computational Materials Science}, vol. 68, pp. 314--319, 2013.
\bibitem{oer} I. C. Man et al., ``Universality in oxygen evolution electrocatalysis on oxide surfaces,'' \emph{ChemCatChem}, vol. 3, pp. 1159--1165, 2011.
\bibitem{her} J. K. Norskov et al., ``Trends in the exchange current for hydrogen evolution,'' \emph{Journal of The Electrochemical Society}, vol. 152, J23, 2005.
\end{thebibliography}

\end{document}
